\hspace{3mm}

\centerline{\bf \Huge Домашнее задание.}
\centerline{\bf \Huge Степень вхождения. Множества.}

\begin{flushright}
    \scriptsize{Настоящий мужчина не умрет, даже если его убить.\\Камина. Гуррен Лаганн.}
\end{flushright}


\problem{1} Натуральные числа $a$ и $b$ таковы, что сумма $\dfrac{a^2}{b} + \dfrac{b^2}{a}$ целая. Докажите, что оба слагаемых целые.

\problem{2} Натуральные числа $a, b, c$ таковы, что число $\dfrac{a}{b} + \dfrac{b}{c} +\dfrac{c}{a}$ является целым. Верно ли, что  $abc \ - $ точный куб?


\problem{3} Найдите все такие составные числа $n$, что для любого разложения на два натуральных сомножителя $n= xy$   сумма $x+y$   является степенью двойки.

\problem{4} Натуральные числа $a, b \ (a<1000)$ таковы, что $a^{21}$ делится на $b^{10}$. Докажите, что $a^2$ делится на $b$.

\problem{5} Докажите, что многочленов с рациональными коэффициентами счётное количество.

\problem{6} Докажите, что множество бесконечных последовательностей состоящих из нулей и единиц не может быть счётным.


\hspace{3mm}
